\documentclass[11pt,a4paper]{article}
% Préambule commun aux documents de PythonIPM/documentation.
% Compilation : tectonic <fichier>.tex
% tectonic compile avec XeTeX, qui lit l'UTF-8 nativement : surtout PAS inputenc ni fontenc,
% faits pour pdfLaTeX. Ce sont eux qui décodaient mal les caractères multi-octets — « cœur »
% ressortait en « cur », et les tirets cadratins en caractères parasites.
\usepackage{lmodern}
\usepackage[french]{babel}
\usepackage{microtype}
\usepackage[a4paper,margin=2.4cm]{geometry}
\usepackage{amsmath,amssymb}
\usepackage{booktabs}
\usepackage{longtable}
\usepackage{array}
\usepackage{ragged2e}
\usepackage{xcolor}
\usepackage{tcolorbox}
\usepackage{enumitem}
\usepackage{fancyhdr}
\usepackage[hidelinks]{hyperref}

\definecolor{encadre}{HTML}{F4F6F8}
\definecolor{bordure}{HTML}{4A6FA5}
\definecolor{alerte}{HTML}{B03A2E}

% Encadré « à retenir » : la lecture non technique du passage qui précède
\newtcolorbox{retenir}[1][]{
  colback=encadre, colframe=bordure, boxrule=0.6pt, arc=2pt,
  left=8pt, right=8pt, top=6pt, bottom=6pt,
  fonttitle=\bfseries, title={#1}
}

% Encadré d'avertissement : un point ouvert, un écart assumé
\newtcolorbox{attention}[1][]{
  colback=white, colframe=alerte, boxrule=0.6pt, arc=2pt,
  left=8pt, right=8pt, top=6pt, bottom=6pt,
  % le bandeau de titre a déjà la couleur d'alerte en fond : un titre `coltitle=alerte`
  % écrirait du rouge sur du rouge. On garde le blanc par défaut, comme l'encadré « retenir ».
  fonttitle=\bfseries, title={#1}
}

\newcommand{\col}[1]{\texttt{#1}}
% \ensuremath et non \textsubscript : \ipm est employé aussi bien dans le texte que dans les
% formules, et « M\textsubscript{0} » placé en mode mathématique est une erreur — silencieuse en
% PDF, mais qui fait échouer la conversion vers Word.
\newcommand{\ipm}{\ensuremath{M_{0}}}

\setlist[itemize]{itemsep=2pt, topsep=4pt}
\setlist[description]{style=nextline, leftmargin=1.2cm, itemsep=5pt}

% La source du document apparaît en tête de chaque page. Les documents RGPH la redéfinissent
% par \renewcommand{\source}{...} AVANT \begin{document}.
\newcommand{\source}{EHCVM 2021}

\pagestyle{fancy}
\fancyhf{}
\fancyhead[L]{\small\itshape IPM Côte d'Ivoire --- \source}
\fancyhead[R]{\small\thepage}
\renewcommand{\headrulewidth}{0.3pt}

\renewcommand{\arraystretch}{1.15}


\renewcommand{\source}{EHCVM 2021 --- variantes Emploi}
\setcounter{tocdepth}{2}

\title{\bfseries Comprendre le calcul de l'IPM pas à pas\\[5pt]
\large Chômage BIT, SU3, agriculture de subsistance et NEET approché}
\author{Pipeline \col{PythonIPM} --- IPM Côte d'Ivoire}
\date{20 septembre 2026}

\begin{document}
\maketitle

\begin{abstract}
\noindent Cette note explique, sans supposer de connaissance préalable, comment les réponses de
l'EHCVM 2021 deviennent des indicateurs de privation au niveau du ménage, comment ces indicateurs
sont pondérés, puis comment sont calculés l'incidence $H$, l'intensité $A$ et l'indice de pauvreté
multidimensionnelle $M_0$. Elle détaille ensuite les trois variantes demandées pour la dimension
Emploi : \textbf{SU3 + agriculture de subsistance}, \textbf{SU3 uniquement} et
\textbf{SU3 + NEET approché}. Les chiffres ont été recalculés sur 12\,965 ménages représentant
29,81 millions de personnes et contrôlés par deux chemins de calcul indépendants.
\end{abstract}

\tableofcontents
\newpage

\section{Ce qu'il faut retenir avant d'entrer dans le détail}

L'IPM ne compte pas simplement le nombre de personnes privées. Il procède en quatre opérations :

\begin{enumerate}
  \item les réponses individuelles et celles du logement sont ramenées au niveau du ménage ;
  \item chaque situation devient un indicateur binaire : 1 si le ménage est privé, 0 sinon ;
  \item les indicateurs sont pondérés, puis additionnés pour obtenir un score entre 0 et 1 ;
  \item le ménage est pauvre si son score atteint au moins un tiers, soit $k=1/3$.
\end{enumerate}

Deux agrégats décrivent ensuite la pauvreté : $H$ est la part de population pauvre et $A$ la part
moyenne de privations subie par les pauvres. Leur produit est l'IPM :

\[
  \boxed{M_0 = H \times A}
\]

\begin{retenir}[Résultat central]
Parmi les trois variantes demandées, \textbf{SU3 + agriculture de subsistance} est la plus
défendable. Elle donne $H=59{,}3~\%$, $A=0{,}4970$ et $M_0=0{,}2949$. Retenir SU3 seul fait
tomber $M_0$ à 0,2438 surtout parce que l'on supprime une privation agricole très fréquente.
La variante SU3 + NEET approché donne 0,2740, mais l'EHCVM ne mesure pas la formation en cours :
l'indicateur construit est un NEE, pas un véritable NEET.
\end{retenir}

\section{Données, unités et pondération de la population}

\subsection{Les bases utilisées}

Quatre bases de l'EHCVM 2021 sont mobilisées :

\begin{center}
\small
\begin{tabular}{@{}p{4.2cm}rp{6.4cm}@{}}
\toprule
\textbf{Base} & \textbf{Lignes} & \textbf{Rôle dans l'IPM} \\
\midrule
Base Individus & 64\,491 & âge, éducation, emploi, état civil, assurance et recours aux soins \\
Base Ménage & 12\,965 & logement, eau, énergie, sanitaires, pièces, poids et géographie \\
Base Avoirs du ménage & 583\,425 & possession des biens d'équipement \\
Base Sécurité alimentaire & 12\,965 utiles & huit questions de l'échelle FIES \\
\bottomrule
\end{tabular}
\end{center}

La clé de fusion est \col{grappe + menage + vague}. Après nettoyage des lignes vides de la base
alimentaire, les quatre sources couvrent les mêmes 12\,965 ménages.

\subsection{Ménage de construction, individu de comptage}

La privation est construite au niveau du \textbf{ménage}. Par exemple, la présence d'un enfant non
scolarisé prive l'ensemble du ménage. Mais les résultats nationaux portent sur les
\textbf{personnes}. Le poids de population du ménage $i$ est :

\[
  n_i = \text{poids d'enquête du ménage}_i \times \text{taille du ménage}_i.
\]

Un ménage de huit personnes représente donc, à poids d'enquête égal, quatre fois plus de
population qu'un ménage de deux personnes. C'est pourquoi le taux de privation ``ménages'' et le
taux pondéré ``population'' peuvent différer.

\begin{attention}[Ne pas confondre deux pourcentages]
L'agriculture de subsistance concerne 48,1~\% des ménages de l'échantillon, mais 40,5~\% de la
population pondérée. Les ménages concernés sont en moyenne d'une taille et d'un poids différents
des autres ménages. Les indices $H$, $A$ et $M_0$ utilisent toujours la pondération population.
\end{attention}

\section{Des réponses brutes aux indicateurs de privation}

Pour chaque indicateur $j$, le pipeline construit :

\[
g^0_{ij} =
\begin{cases}
1 & \text{si le ménage } i \text{ est privé selon le seuil } z_j,\\
0 & \text{sinon.}
\end{cases}
\]

Un ménage sans personne appartenant à la tranche d'âge d'un indicateur est classé non privé sur
cet indicateur. Ainsi, un ménage sans membre de 17 à 40 ans n'est pas privé au titre de SU3.

\subsection{Les quatre indicateurs d'éducation}

\begin{description}
  \item[Fréquentation scolaire.] Un ménage est privé si au moins un enfant de 6 à 16 ans ne
  fréquente pas l'école. La scolarisation combine la question 2.08a pour 2021/22 et la question
  2.12 pour 2020/21. Privation : 27,9~\% des ménages, 31,9~\% de la population.

  \item[Années de scolarité.] Les années accomplies sont reconstruites à partir du niveau et de
  la classe, achevés ou en cours. Le ménage est privé si aucun membre de 17 à 95 ans n'a accompli
  dix années d'études. Privation : 72,5~\% des ménages, 62,9~\% de la population.

  \item[Alphabétisation.] Est alphabétisé celui qui sait à la fois lire et écrire le français.
  Le ménage est privé si au moins un membre de 17 à 49 ans ne satisfait pas ces deux conditions.
  Privation : 61,7~\% des ménages, 62,9~\% de la population.

  \item[État civil.] Le ménage est privé si au moins un enfant de 5 à 15 ans ne possède pas
  d'acte de naissance. Privation : 23,5~\% des ménages, 26,3~\% de la population.
\end{description}

\subsection{Les trois indicateurs de santé}

\begin{description}
  \item[Assurance maladie.] Le ménage est privé si aucun membre n'est couvert. Privation :
  93,4~\% des ménages, 90,2~\% de la population.

  \item[Insécurité alimentaire.] Les huit réponses FIES sont additionnées. Un score d'au moins
  quatre ``oui'' signale une insécurité modérée ou sévère. Privation : 42,7~\% des ménages,
  40,8~\% de la population.

  \item[Renoncement aux soins.] Il faut qu'un membre ait été malade au cours des 30 derniers
  jours, n'ait pas consulté et cite une raison subie : soins trop chers, manque d'argent, service
  indisponible ou absence de personnel. Privation : 9,3~\% des ménages, 9,0~\% de la population.
\end{description}

\subsection{Les sept indicateurs de conditions de vie}

\begin{description}
  \item[Électricité.] Privé si l'éclairage principal ne vient ni du réseau, ni d'un groupe
  électrogène, ni du solaire. Population privée : 9,3~\%.

  \item[Logement.] Privé si le sol est naturel ou si le toit ou les murs sont constitués de
  matériaux naturels ou rudimentaires. Population privée : 17,3~\%.

  \item[Eau potable.] Privé si la source n'est pas améliorée selon la classification ODD, ou si
  elle se trouve à plus de 15 minutes à l'aller, soit plus de 30 minutes aller-retour.
  Population privée : 20,9~\%.

  \item[Énergie de cuisson.] Privé si le combustible principal n'est ni le gaz ni l'électricité.
  Population privée : 66,3~\%.

  \item[Toilettes.] Privé si les sanitaires ne sont pas améliorés ou s'ils sont partagés avec
  d'autres ménages. Population privée : 67,2~\%.

  \item[Biens d'équipement.] Privé si le ménage possède au plus un bien parmi radio, télévision,
  téléphone, ordinateur, vélo, moto et réfrigérateur, et ne possède pas de voiture. La charrette,
  prévue par la définition PNUD, n'existe pas dans le module EHCVM. Population privée : 21,7~\%.

  \item[Promiscuité.] Privé si le nombre de personnes par pièce à coucher dépasse trois.
  Population privée : 12,5~\%.
\end{description}

\section{Les indicateurs d'emploi, dans le détail}

\subsection{Chômage BIT ou SU1}

Une personne est chômeuse si trois conditions sont simultanément réunies :

\begin{enumerate}
  \item \textbf{sans emploi} : elle n'a pas travaillé au moins une heure pendant les sept
  derniers jours et n'a pas un emploi dont elle était temporairement absente ;
  \item \textbf{en recherche} : elle a recherché un emploi rémunéré dans les 30 derniers jours ;
  \item \textbf{disponible} : elle peut commencer immédiatement, sous 15 jours ou entre 15 jours
  et un mois, ou se déclare disponible immédiatement dans la branche des non-chercheurs.
\end{enumerate}

Pour l'IPM, seules les personnes de 17 à 40 ans sont retenues. Le ménage est privé dès qu'il en
compte au moins une. On obtient 561 chômeurs dans cette tranche, 476 ménages privés et 6,5~\% de
population privée.

Le chômage BIT est rare car une seule heure de travail suffit à être classé occupé et parce que
les personnes découragées qui ne recherchent plus activement sortent du chômage.

\subsection{SU3 : chômage plus main-d'œuvre potentielle}

SU3 ajoute au chômage la \textbf{main-d'œuvre potentielle}. Celle-ci comprend :

\begin{itemize}
  \item les personnes sans emploi qui cherchent mais ne sont pas disponibles ;
  \item les personnes sans emploi qui sont disponibles mais ne cherchent pas, à condition
  qu'elles ne déclarent pas explicitement ne pas vouloir travailler.
\end{itemize}

Au niveau individuel :

\[
  \text{SU3} = \text{chômage BIT} \ \cup\ \text{main-d'œuvre potentielle}.
\]

Tout chômeur BIT appartient donc nécessairement à SU3. Cette propriété a été vérifiée dans les
données : aucune des 476 privations de ménage au titre du chômage ne se trouve hors de SU3.

Sur les 17--40 ans, SU3 concerne 1\,200 personnes, contre 561 pour le chômage BIT. À l'échelle du
ménage, 1\,027 ménages sont privés, soit 7,9~\% des ménages et 12,1~\% de la population.

\begin{attention}[Pourquoi il ne faut pas additionner mécaniquement SU1 et SU3]
SU3 contient déjà SU1. Une variante ``SU1 + SU3'' attribue donc deux privations au même chômage
pour les ménages concernés. Ce montage peut être calculé comme test, mais il n'est pas recommandé
comme spécification principale.
\end{attention}

\subsection{Agriculture de subsistance : ce que mesure réellement le code}

Une personne est classée en agriculture de subsistance si elle :

\begin{itemize}
  \item a travaillé pendant les sept derniers jours ;
  \item a travaillé dans son propre champ ;
  \item n'a exercé ni activité salariée, ni apprentissage, ni commerce.
\end{itemize}

Le ménage n'est privé que si cette personne est le \textbf{chef de ménage}. Cette règle donne
6\,238 ménages privés, soit 48,1~\% des ménages et 40,5~\% de la population.

Le terme ``subsistance'' doit être manié avec prudence. La définition statistique de l'OIT repose
sur la destination principale de la production : consommation du ménage plutôt que marché.
L'EHCVM observe le travail exclusif dans son propre champ, mais pas cette destination. Le code
mesure donc plus exactement une \textbf{dépendance exclusive du chef à l'agriculture sur son
propre champ}.

\subsection{NEET approché : en réalité un NEE}

L'indicateur international NEET porte sur les jeunes de 15 à 24 ans qui ne sont :

\begin{itemize}
  \item ni en emploi ;
  \item ni en éducation ;
  \item ni en formation.
\end{itemize}

L'emploi et la scolarisation sont observables dans l'EHCVM. En revanche, les questions de
formation professionnelle demandent si la personne en a déjà suivi une, sans établir qu'elle en
suit une actuellement. Le pipeline construit donc :

\[
  \text{NEE approché} = (15\text{--}24\text{ ans}) \cap
  (\text{sans emploi}) \cap (\text{non scolarisé}).
\]

On obtient 4\,042 jeunes sur 10\,078, soit 40,1~\%. Au niveau du ménage, 3\,247 ménages sont
concernés, représentant 29,1~\% de la population. Le chiffre est une borne supérieure du NEET :
des jeunes comptés NEE peuvent suivre une formation non observée.

Par ailleurs, SU3 et NEE se recouvrent dans 547 ménages. Le couple SU3 + NEE mesure deux objets
différents, mais donne deux fois du poids à certains épisodes de non-emploi des jeunes.

\section{Pondérations : comment chaque privation entre dans le score}

Les quatre dimensions sont équipondérées : chacune reçoit 0,25. À l'intérieur d'une dimension,
ce poids est partagé également entre les indicateurs retenus.

\subsection{IPM national de référence}

\begin{center}
\small
\begin{tabular}{@{}lrrr@{}}
\toprule
\textbf{Dimension} & \textbf{Nombre d'indicateurs} & \textbf{Poids par indicateur} &
\textbf{Poids total} \\
\midrule
Éducation & 4 & 0,0625 & 0,25 \\
Santé & 3 & 0,083333 & 0,25 \\
Emploi & 2 & 0,1250 & 0,25 \\
Conditions de vie & 7 & 0,035714 & 0,25 \\
\midrule
Total & 16 & --- & 1,00 \\
\bottomrule
\end{tabular}
\end{center}

La somme des poids vaut exactement 1. Dans la référence, les deux indicateurs d'emploi sont le
chômage BIT et l'agriculture de subsistance.

\subsection{Poids dans les trois variantes demandées}

\begin{center}
\small
\begin{tabular}{@{}p{4.3cm}p{7cm}p{2.1cm}@{}}
\toprule
\textbf{Variante} & \textbf{Indicateurs d'emploi retenus} & \textbf{Poids} \\
\midrule
SU3 + agriculture & SU3 ; agriculture sur propre champ exclusive du chef & 0,125 chacun \\
SU3 uniquement & SU3 & 0,250 \\
SU3 + NEET approché & SU3 ; NEE des 15--24 ans & 0,125 chacun \\
\bottomrule
\end{tabular}
\end{center}

Les poids des trois autres dimensions ne changent jamais. Seul le partage des 0,25 de la
dimension Emploi varie.

\section{Du score du ménage à la pauvreté multidimensionnelle}

\subsection{Score pondéré}

Pour le ménage $i$, le score est :

\[
  c_i = \sum_{j=1}^{J} w_j g^0_{ij},
\]

où $w_j$ est le poids de l'indicateur et $g^0_{ij}$ vaut 0 ou 1. Un score de 0,40 signifie que le
ménage subit 40~\% de l'ensemble pondéré des privations possibles. Il ne signifie pas que 40~\%
des membres du ménage sont privés.

\subsection{Seuil de pauvreté}

Le seuil retenu est $k=1/3$. Le statut de pauvreté est :

\[
  \rho_i(k) = \mathbf{1}(c_i \geq 1/3).
\]

Un score exactement égal à un tiers est donc pauvre. Deux catégories complémentaires sont aussi
calculées : vulnérable si $0,20 \leq c_i < 1/3$ et pauvreté sévère si $c_i \geq 0,50$.

\subsection{Censure des privations}

Pour expliquer la composition de l'IPM, les privations des ménages non pauvres sont remises à
zéro :

\[
  g^0_{ij}(k) = g^0_{ij}\rho_i(k).
\]

C'est la censure Alkire--Foster. Une privation réelle d'un ménage non pauvre compte dans le taux
brut de l'indicateur, mais pas dans $M_0$ ni dans les contributions à la pauvreté
multidimensionnelle.

\section{Calcul de H, A et M0}

Soit $n_i$ le poids de population du ménage et $q_i=\rho_i(k)$ son statut de pauvreté.

\subsection{Incidence H}

\[
  H = \frac{\sum_i n_i q_i}{\sum_i n_i}.
\]

$H$ répond à la question : quelle part de la population vit dans un ménage dont le score atteint
au moins un tiers ?

\subsection{Intensité A}

\[
  A = \frac{\sum_i n_i c_i q_i}{\sum_i n_i q_i}.
\]

$A$ est le score moyen des seuls pauvres. Une intensité de 0,497 signifie que les personnes
pauvres subissent en moyenne 49,7~\% des privations pondérées possibles.

\subsection{Indice M0}

\[
  M_0 = H A = \frac{\sum_i n_i c_i q_i}{\sum_i n_i}.
\]

$M_0$ augmente si davantage de personnes deviennent pauvres, si les pauvres accumulent davantage
de privations, ou si les deux phénomènes se produisent ensemble.

\begin{attention}[M0 n'est pas un pourcentage ordinaire]
Un $M_0$ de 0,2949 ne signifie pas que 29,49~\% des personnes sont pauvres. Ici, 59,3~\% sont
pauvres et leur intensité moyenne est de 49,70~\%. Le produit $0,5932\times0,4970$ donne 0,2949.
\end{attention}

\section{Exemple entièrement calculé pour un ménage fictif}

Considérons la variante SU3 + agriculture. Supposons un ménage privé sur six indicateurs : année
de scolarité, alphabétisation, assurance maladie, SU3, énergie de cuisson et toilettes.

\begin{center}
\small
\begin{tabular}{@{}lccc@{}}
\toprule
\textbf{Indicateur privé} & $g^0_{ij}$ & $w_j$ & $w_jg^0_{ij}$ \\
\midrule
Année de scolarité & 1 & 0,062500 & 0,062500 \\
Alphabétisation & 1 & 0,062500 & 0,062500 \\
Assurance maladie & 1 & 0,083333 & 0,083333 \\
SU3 & 1 & 0,125000 & 0,125000 \\
Énergie de cuisson & 1 & 0,035714 & 0,035714 \\
Toilettes & 1 & 0,035714 & 0,035714 \\
\midrule
Score $c_i$ & & & \textbf{0,404761} \\
\bottomrule
\end{tabular}
\end{center}

Comme $0,404761 \geq 1/3$, le ménage est pauvre. Son score censuré reste 0,404761. Si son score
avait été 0,30, il aurait été vulnérable mais non pauvre et toutes ses privations auraient été
remises à zéro dans le calcul de $M_0$.

Dans la variante SU3 uniquement, la privation SU3 vaudrait 0,25 au lieu de 0,125. Toutes choses
égales par ailleurs, le score fictif passerait à 0,529761. Cela montre pourquoi une modification
du nombre d'indicateurs change non seulement la description de l'emploi, mais aussi la position de
certains ménages par rapport aux seuils d'un tiers et de moitié.

\section{Résultats nationaux des variantes}

\begin{center}
\small
\begin{tabular}{@{}p{4.8cm}rrrr@{}}
\toprule
\textbf{Dimension Emploi} & $H$ & $A$ & $M_0$ & \textbf{Sévère} \\
\midrule
Référence : chômage BIT + agriculture & 58,0~\% & 0,4940 & 0,2863 & 26,1~\% \\
SU3 + agriculture de subsistance & 59,3~\% & 0,4970 & \textbf{0,2949} & 27,1~\% \\
SU3 uniquement & 53,6~\% & 0,4546 & \textbf{0,2438} & 13,7~\% \\
SU3 + NEET approché (NEE) & 57,7~\% & 0,4746 & \textbf{0,2740} & 20,5~\% \\
\bottomrule
\end{tabular}
\end{center}

Les intervalles de confiance à 95~\% de $M_0$ sont respectivement :

\begin{itemize}
  \item référence : $[0,2701\ ;\ 0,3026]$ ;
  \item SU3 + agriculture : $[0,2788\ ;\ 0,3109]$ ;
  \item SU3 uniquement : $[0,2331\ ;\ 0,2545]$ ;
  \item SU3 + NEE : $[0,2617\ ;\ 0,2863]$.
\end{itemize}

\subsection{Pourquoi SU3 + agriculture augmente légèrement l'IPM}

SU3 remplace le chômage BIT au même poids de 0,125. La privation population passe de 6,5~\% à
12,1~\%. Certains ménages auparavant juste sous un tiers franchissent le seuil ; $H$ passe de
58,0~\% à 59,3~\%. Chez les pauvres, SU3 ajoute aussi un peu de privation, ce qui porte $A$ de
0,4940 à 0,4970.

\subsection{Pourquoi SU3 uniquement fait fortement baisser l'IPM}

SU3 reçoit 0,25, mais la privation agricole, qui concernait 40,5~\% de la population, disparaît.
La hausse du poids de SU3 ne compense pas cette suppression. Des ménages ruraux perdent 0,125 de
score et repassent sous le seuil. La pauvreté sévère est particulièrement sensible : elle passe
de 27,1~\% dans SU3 + agriculture à 13,7~\% dans SU3 uniquement.

Cette baisse est un effet de définition et de pondération. Elle ne représente aucune amélioration
des conditions de vie entre deux dates.

\subsection{Pourquoi SU3 + NEE donne un résultat intermédiaire}

Le NEE est beaucoup plus fréquent que SU3 : sa privation brute touche 29,1~\% de la population au
niveau ménage. Il remplace donc une partie du pouvoir discriminant de l'agriculture, mais déplace
la dimension vers la jeunesse et la déconnexion de l'école. La variante devient moins rurale et
plus sensible aux ménages comprenant des jeunes de 15 à 24 ans.

\section{Lecture urbaine et rurale}

\begin{center}
\small
\begin{tabular}{@{}p{4.6cm}rrrr@{}}
\toprule
& \multicolumn{2}{c}{\textbf{Urbain}} & \multicolumn{2}{c}{\textbf{Rural}} \\
\cmidrule(lr){2-3}\cmidrule(l){4-5}
\textbf{Variante} & $H$ & $M_0$ & $H$ & $M_0$ \\
\midrule
Référence & 33,6~\% & 0,1525 & 85,1~\% & 0,4354 \\
SU3 + agriculture & 35,7~\% & 0,1626 & 85,6~\% & 0,4422 \\
SU3 uniquement & 37,6~\% & 0,1685 & 71,5~\% & 0,3278 \\
SU3 + NEE & 40,9~\% & 0,1850 & 76,5~\% & 0,3732 \\
\bottomrule
\end{tabular}
\end{center}

La suppression de l'agriculture réduit surtout l'IPM rural. À l'inverse, NEE fait davantage
remonter l'incidence urbaine. Les variantes ne changent donc pas seulement le niveau national :
elles modifient aussi le profil territorial de la pauvreté et, par conséquent, le ciblage public.

\section{Contribution des indicateurs à M0}

La contribution de l'indicateur $j$ s'écrit :

\[
  \mathrm{Contribution}_j =
  \frac{w_j\,CH_j}{M_0},
\]

où $CH_j$ est le taux de privation censuré : la part de population qui est à la fois pauvre et
privée sur l'indicateur $j$. La somme des contributions vaut 100~\%.

\begin{center}
\small
\begin{tabular}{@{}lrrrr@{}}
\toprule
\textbf{Contribution à M0} & \textbf{Référence} & \textbf{SU3+agri.} &
\textbf{SU3 seul} & \textbf{SU3+NEE} \\
\midrule
Chômage BIT & 1,6~\% & --- & --- & --- \\
SU3 & --- & 3,2~\% & 11,2~\% & 4,0~\% \\
Agriculture de subsistance & 16,9~\% & 16,5~\% & --- & --- \\
NEET approché (NEE) & --- & --- & --- & 11,9~\% \\
\midrule
Dimension Emploi & 18,5~\% & 19,7~\% & 11,2~\% & 15,9~\% \\
\bottomrule
\end{tabular}
\end{center}

Le poids nominal de la dimension reste 25~\%, mais sa contribution observée à $M_0$ dépend de la
fréquence de ses privations et de leur chevauchement avec les autres dimensions.

\section{Tableau de contrôle de l'IPM national de référence}

Le tableau suivant relie règle, fréquence, poids et contribution. Le taux population est non
censuré ; la contribution ne porte que sur les pauvres après censure.

\begin{center}
\scriptsize
\begin{longtable}{@{}p{3.4cm}p{2.6cm}rrr@{}}
\toprule
\textbf{Indicateur} & \textbf{Dimension} & \textbf{Privation pop.} & $w_j$ &
\textbf{Contribution} \\
\midrule
\endfirsthead
\toprule
\textbf{Indicateur} & \textbf{Dimension} & \textbf{Privation pop.} & $w_j$ &
\textbf{Contribution} \\
\midrule
\endhead
Fréquentation scolaire & Éducation & 31,9~\% & 0,0625 & 6,0~\% \\
Années de scolarité & Éducation & 62,9~\% & 0,0625 & 10,8~\% \\
Alphabétisation & Éducation & 62,9~\% & 0,0625 & 10,3~\% \\
État civil & Éducation & 26,3~\% & 0,0625 & 5,2~\% \\
Assurance maladie & Santé & 90,2~\% & 0,0833 & 16,7~\% \\
Insécurité alimentaire & Santé & 40,8~\% & 0,0833 & 9,0~\% \\
Renoncement aux soins & Santé & 9,0~\% & 0,0833 & 2,2~\% \\
Chômage BIT & Emploi & 6,5~\% & 0,1250 & 1,6~\% \\
Agriculture de subsistance & Emploi & 40,5~\% & 0,1250 & 16,9~\% \\
Électricité & Conditions de vie & 9,3~\% & 0,0357 & 1,1~\% \\
Logement & Conditions de vie & 17,3~\% & 0,0357 & 2,0~\% \\
Eau potable & Conditions de vie & 20,9~\% & 0,0357 & 2,2~\% \\
Énergie de cuisson & Conditions de vie & 66,3~\% & 0,0357 & 6,5~\% \\
Toilettes & Conditions de vie & 67,2~\% & 0,0357 & 6,3~\% \\
Biens d'équipement & Conditions de vie & 21,7~\% & 0,0357 & 2,1~\% \\
Promiscuité & Conditions de vie & 12,5~\% & 0,0357 & 1,1~\% \\
\midrule
Total des contributions & & & & 100,0~\% \\
\bottomrule
\end{longtable}
\end{center}

\section{La littérature soutient-elle l'indicateur agricole ?}

\subsection{Arguments favorables}

L'OIT souligne que le chômage ouvert peut être peu informatif dans les économies où les personnes
pauvres ne peuvent pas se permettre de rester sans aucune activité. Elle propose, parmi les
indicateurs adaptés aux moyens d'existence agricoles, la part des ménages engagés dans
l'agriculture de subsistance. Elle insiste aussi sur la nécessité d'adapter les indicateurs à la
structure nationale de l'emploi et des moyens d'existence.\footnote{OIT, \emph{Jobs and
livelihoods: Meaningful ways to set targets and monitor progress},
\url{https://www.ilo.org/resource/jobs-and-livelihoods-meaningful-ways-set-targets-and-monitor-progress}.}

Le Rwanda offre un précédent officiel d'IPM national. Son indicateur de subsistance prive le
ménage lorsqu'il est uniquement engagé dans le travail agricole. Il appartient à la dimension
``Services sociaux et activité économique''.\footnote{National Institute of Statistics of
Rwanda, \emph{The Second Rwanda Multidimensional Poverty Report}, p.~31,
\url{https://ophi.org.uk/sites/default/files/2023-12/nisr_rwanda_mpi_report_2018.pdf}.}

Enfin, le guide PNUD--OPHI rappelle qu'un IPM national peut inclure l'emploi et d'autres dimensions
adaptées aux priorités du pays.\footnote{PNUD--OPHI, \emph{How to build a National
Multidimensional Poverty Index},
\url{https://www.undp.org/publications/how-build-national-multidimensional-poverty-index}.}

\subsection{Réserve de définition}

Depuis la 19e Conférence internationale des statisticiens du travail, l'emploi désigne le travail
pour rémunération ou profit. La production principalement destinée à l'usage propre est une autre
forme de travail, à mesurer séparément.\footnote{OIT, \emph{Employment},
\url{https://www.ilo.org/resource/employment-1} ; OIT, \emph{Own-use production work},
\url{https://www.ilo.org/resource/own-use-production-work}.}

La littérature soutient donc clairement l'idée de mesurer la dépendance à la production agricole
de subsistance. Elle ne suffit pas à valider le proxy EHCVM, car celui-ci ne connaît pas la
destination de la récolte. Le libellé conseillé est :

\begin{retenir}[Libellé recommandé]
« Chef de ménage dépendant exclusivement de l'agriculture sur son propre champ --- proxy de
l'agriculture de subsistance. »
\end{retenir}

\section{Forces, limites et risques d'interprétation}

\begin{description}
  \item[SU3 est préférable au chômage BIT strict.] Il récupère une partie du découragement et de
  l'attachement incomplet au marché du travail. Il reste toutefois limité aux 17--40 ans et la
  disponibilité admise jusqu'à un mois est plus large que le standard de deux semaines.

  \item[L'agriculture mesure une vulnérabilité absente de SU3.] Une personne travaillant dans son
  champ est occupée et n'entre pas dans SU3. Sans indicateur de qualité de l'emploi, cette réalité
  disparaît de la dimension.

  \item[SU3 seul donne un poids très fort à une privation rare.] Chaque ménage SU3 reçoit 0,25 de
  score d'un coup. Cela rend l'indice sensible à une seule condition tout en ignorant la qualité
  des emplois existants.

  \item[Le NEE n'est pas un NEET complet.] La formation en cours manque et le taux est donc
  surestimé. De plus, l'indicateur ne concerne que les ménages comprenant un jeune de 15 à 24 ans.

  \item[Les variantes changent le profil des pauvres.] Supprimer l'agriculture réduit surtout la
  pauvreté rurale ; ajouter NEE augmente relativement le poids des ménages jeunes et urbains.

  \item[Une comparaison entre variantes n'est pas une évolution temporelle.] Les ménages et la
  date restent identiques. Seule la définition de l'indice change.
\end{description}

\section{Recommandation finale}

\begin{enumerate}
  \item Retenir comme variante principale \textbf{SU3 + dépendance agricole exclusive du chef},
  avec deux indicateurs de poids 0,125.
  \item Renommer ou qualifier explicitement l'indicateur agricole comme \textbf{proxy}, car la
  destination de la production n'est pas observée.
  \item Conserver \textbf{SU3 uniquement} comme test de sensibilité montrant ce qui arrive lorsque
  la qualité de l'emploi agricole est retirée.
  \item Conserver \textbf{SU3 + NEE} comme scénario exploratoire, sans l'appeler NEET dans une
  publication officielle tant que la formation en cours n'est pas mesurée.
  \item À la prochaine enquête, ajouter une question sur la destination principale de la
  production agricole et une question sur la formation suivie actuellement. Ces deux ajouts
  permettraient de construire les indicateurs sans approximation.
\end{enumerate}

\section{Contrôles réalisés pour sécuriser les chiffres}

Les vérifications suivantes ont été exécutées après l'implémentation :

\begin{itemize}
  \item auto-contrôles unitaires des cinq étapes du pipeline, sur l'EHCVM et le RGPH ;
  \item recalcul direct et séparé des scores, de $H$, de $A$ et de $M_0$ depuis la matrice binaire ;
  \item vérification exacte de l'identité $M_0=H\times A$ ;
  \item écart maximal de $5,15\times10^{-7}$ entre le recalcul indépendant et les CSV publiés,
  entièrement expliqué par leur arrondi à six décimales ;
  \item contrôle que tout ménage privé au chômage BIT est également privé selon SU3 ;
  \item validation de l'IPM national avec le package indépendant \col{afmpi}, y compris les
  contributions et l'erreur-type, avec une tolérance de $10^{-6}$ ;
  \item vérification que SU3 et NEE restent hors de l'IPM national lorsqu'aucune option de
  variante n'est demandée.
\end{itemize}

\begin{retenir}[Conclusion]
Les valeurs présentées dans cette note sont reproductibles et cohérentes avec les règles codées.
Le choix restant n'est pas informatique : il consiste à décider quelle conception de la
privation d'emploi doit entrer dans l'IPM. Les résultats et la littérature conduisent à préférer
SU3 accompagné d'un indicateur agricole, sous un libellé plus prudent et statistiquement exact.
\end{retenir}

\section*{Fichiers de résultats associés}

\begin{itemize}
  \item \col{indices\_ipm\_su3\_ehcvm2021.xlsx} : SU3 + agriculture ;
  \item \col{indices\_ipm\_su3\_seul\_ehcvm2021.xlsx} : SU3 uniquement ;
  \item \col{indices\_ipm\_su3\_neet\_ehcvm2021.xlsx} : SU3 + NEE ;
  \item chaque classeur contient les résultats nationaux, les contributions, les poids et les
  désagrégations territoriales et sociales.
\end{itemize}

\end{document}
